\section{Plans, forecasts and repeated errors}
\label{sec:temporal}

The record also contains what the systems said they would do, what they expected of others, and how they responded to rejected actions. This evidence is thinner.

\paragraph{Stated plans meet opportunities unevenly.} An earlier pass read the public message and private note of 240 sampled decisions, three per player per game. A language-model reader marked forward-looking statements about the player's own actions in 59. Most are contingent or open-ended. Eight state a checkable policy of buying a named property on landing there. Following each player to the end of its game, six never get the chance. One is carried out. The eighth meets its opportunity and is not. In game 16 Grok's note includes the plan \InlineQuote{fx-g16-plan}. Fifty-six decisions later it lands on Virginia Avenue with \$371 against a \$160 price. Its note says \InlineQuote{fx-g16-note} Its public message says \InlineQuote{fx-g16-public} The tool call starts an auction, which Gemini wins at \$380. This is one of the eleven Grok decisions whose text buys and whose call auctions. A plan can fail for lack of an opportunity or at the response layer, and neither is a change of mind. Opus's notes are absent in 53 of its 60 sampled decisions, so these counts cannot compare planning across systems.

\paragraph{Some targets are carried across many turns.} In game 9, Sol's note at turn 29 sets a target of three houses on each dark-blue deed as funds arrive. It adds one house the next turn, another at turn 39, and reaches three on each at turn 60 after mortgaging a spare deed. The target is realized through intervening rolls, rents and financing. The chronology does not show that the note caused the result.

\paragraph{Statements about opponents are mostly accounting.} The same reader marked statements about what an opponent could or would do in 36 sampled decisions. Nearly all concern public quantities such as cash, reachable rents and mortgage capacity. Some are checkable and correct. In the Boardwalk purchase of Section~\ref{sec:cashpoor}, Grok's note predicts that Opus cannot survive a light-blue hotel, and four turns later Opus is eliminated owing Gemini \$550. In Incident~\ref{inc:swap}, Sol's third offer answers the objection Grok had stated in public. These show attention to a rival's finances and stated reasons. They are not predictions of hidden beliefs.

\paragraph{An encoding error recurs when feedback is absent, and sometimes when it is present.} The uneven-build rejection of Section~\ref{sec:defaults} occurs in 114 decisions. For each we find the same player's next decision in that game that submits a build plan (Table~\ref{tab:recurrence}). The error recurs at that next submission in 63 of 109 cases. Fifty-four of the 63 are Opus. Its retry had been spent on the refusal, so in 60 of its 61 error decisions no request showed it the validator's message. In the 54 decisions where a corrective request did show the message, the error recurs at 9 of 50 next submissions, and at 2 of 29 when only the first such request in each player-game is counted. These are sequences under changing states with no control condition. The error is usually not repeated once the message has been seen, which is not a measurement of learning.

\begin{table}[t]
\centering\small
\begin{tabular}{lrrr}
\toprule
Uneven-build rejections & Decisions & Later submission & Error repeated\\
\midrule
All & 114 & 109 & 63\\
Validator message shown in a corrective request & 54 & 50 & 9\\
First such request per player and game & 31 & 29 & 2\\
\bottomrule
\end{tabular}
\caption{Recurrence of the uneven-build encoding error at the same player's next build submission in the same game. The populations are nested. Decisions with no later submission are censored. Inherited from the earlier audit.}
\label{tab:recurrence}
\end{table}
